\begin{sddsprog}{sddseditor}
  \item \textbf{description:} \verb|sddseditor| is a Qt graphical application for viewing
  and modifying SDDS files. Multiple pages of parameters, columns and arrays can be
  edited using familiar table widgets. Standard shortcuts such as
  \verb|Ctrl+N| for creating files, \verb|Ctrl+O| for opening files,
  \verb|Ctrl+S| for saving, \verb|Ctrl+C| and
  \verb|Ctrl+V| for copy and paste, and \verb|Ctrl+Z|/\verb|Ctrl+Shift+Z| for
  undo and redo are provided. Additional actions are available from context menus
  on the column and array headers, including sorting columns, plotting with
  \progref{sddsplot}, and searching and resizing arrays. Files may also be
  exported as comma-separated values using \verb|Ctrl+Shift+C| and to HDF format
  using \verb|Ctrl+Shift+H|.

  Choose \verb|File > New...| or the \verb|New| toolbar button to create a file,
  even while another file is open. The setup dialog lists parameters, columns
  and arrays in separate cards, each with names, types and units. Parameters
  take an initial value, columns share a row count, and arrays take a shape such
  as \verb|5| or \verb|3x4|. An optional description is stored in the file
  header. Problems such as duplicate names are reported as you type, and
  \verb|Create| is enabled once every definition is valid. New files start with
  one page. Removing all definitions creates an empty file. Unsaved changes in
  the current document are protected by a Save/Discard/Cancel prompt, and saving
  the new file asks for a filename.

  \verb|Edit > Column Rows > Filter/View...| (\verb|Ctrl+Shift+R|) shows only the
  column rows that match an expression such as \verb|X > 0 && Status == "OK"|. No
  data is deleted, and copy, paste, delete, search/replace and the formula tools
  then change only the visible rows. The \verb|Edit > Formula / Fill| menu sets
  the selected cells to a series (\verb|Ctrl+Shift+F|), the value of a numerical
  expression such as \verb|x * 2| (\verb|Ctrl+Shift+E|), or a text template such
  as \verb|run_${i}| (\verb|Ctrl+Shift+M|). Multidimensional arrays can be browsed
  one slice at a time with \verb|Edit > Array > Open Array Viewer...|. Numeric
  columns whose names contain \verb|Time| and whose units are seconds or blank can
  show their values as local or UTC dates from the \verb|Time display| context
  menu; this changes only the display, not the saved data.

  SDDS6 enum parameters, columns and arrays are shown as \verb|label (code)|, and
  their values are chosen from the declared labels. The element's \verb|Attributes|
  dialog edits its code-to-label mapping, such as \verb|0:OFF,10:ON,-1:FAULT|. The
  editor does not create enum elements or change their type; use
  \progref{sddsmakedataset} and \progref{sddscast} for that, and convert enums with
  \progref{sddscast} before exporting to CSV or HDF. \verb|Info > Help| describes
  these features in more detail.

  \item \textbf{examples:}
    \begin{verbatim}
    sddseditor lattice.twi
    sddseditor
    \end{verbatim}

  \item \textbf{synopsis:}
    \begin{verbatim}
    sddseditor [inputFile]
    \end{verbatim}

  \item \textbf{see also:}
    \begin{itemize}
      \item \progref{sddsplot}
      \item \progref{sddsprocess}
      \item \progref{sddscast}
      \item \progref{sddsmakedataset}
    \end{itemize}
  \item \textbf{author:} R. Soliday, ANL/APS.
\end{sddsprog}
